% These packages are required for list code snippets.
\usepackage{xcolor}
\usepackage{color}
% The following special color definitions are used in the IST Thesis
\definecolor{forestgreen}{RGB}{34,139,34}
\definecolor{orangered}{RGB}{239,134,64}
\definecolor{lightred}{rgb}{1,0.4,0.5}
\definecolor{orange}{rgb}{1,0.45,0.13}
\definecolor{darkblue}{rgb}{0.0,0.0,0.6}
\definecolor{lightblue}{rgb}{0.1,0.57,0.7}
\definecolor{gray}{rgb}{0.4,0.4,0.4}
\definecolor{lightgray}{rgb}{0.95, 0.95, 0.95}
\definecolor{darkgray}{rgb}{0.4, 0.4, 0.4}
\definecolor{editorGray}{rgb}{0.95, 0.95, 0.95}
\definecolor{editorOcher}{rgb}{1, 0.5, 0} % #FF7F00 -> rgb(239, 169, 0)
\definecolor{chaptergrey}{rgb}{0.6,0.6,0.6}
\definecolor{editorGreen}{rgb}{0, 0.5, 0} % #007C00 -> rgb(0, 124, 0)
\definecolor{olive}{rgb}{0.17,0.59,0.20}
\definecolor{brown}{rgb}{0.69,0.31,0.31}
\definecolor{purple}{rgb}{0.38,0.18,0.81}

\newcommand{\changed}[1]{\textcolor{orange}{#1}}

% Symbols for linear logic
\usepackage{cmll}

% For proof trees
\usepackage{bussproofs}

% For multiply counters, which I use for indentation in class environments
\usepackage{calc}

% These are useful!
\providecommand\todo{}
\renewcommand{\todo}[1]{\textcolor{red}{\textbf{#1}}}

% Useful for CLASS snippets
\lstdefinelanguage{CLASS}{
    morekeywords={proc, cut, close, wait, case, of, send, recv, choice, offer, !, ?, print, println, par, let, share, put, take, cell},
    morekeywords=[2]{lint, colint, lbool, colbool, state, usage},
    sensitive=true,
    moredelim=[is][keywordstyle2]{@}{@},
}

% Useful for IR snippets
\lstdefinelanguage{IR}{
    morekeywords={int, bool, string, tag, cont, exponential, cell, type, isValue},
    morecomment=[l]{//},
    sensitive=true,
}

% Color BNF with strings
\lstdefinelanguage{BNF}{
    morekeywords={nl, pid, lab, int, loc, t, t1, t2, c, c1, c2, d, d1, d2, e, prog, proc, entry, block, instr, sl, st, sc, ofs, locb, reqs, flag, expr, tflg, targ, carg, darg, parg, pargs, earg, eargs, case, if, else, str, bool, src, dst},
    sensitive=true,
    morestring=[b]",
    morestring=[b]',
    morecomment=[l]{//},
    morecomment=[s]{/*}{*/},
    moredelim=[is][keywordstyle2]{@}{@},
    alsoletter={:},
    morekeywords=[2]{::=,|,*,+,?},
}

% Slot names
\newcommand{\irs}[1]{\text{\textcolor{lightblue}{#1}}}
\newcommand{\ints}{\irs{int}}
\newcommand{\bools}{\irs{bool}}
\newcommand{\strings}{\irs{string}}
\newcommand{\tags}{\irs{tag}}
\newcommand{\conts}{\irs{cont}}
\newcommand{\exponentials}{\irs{exponential}}
\newcommand{\types}{\irs{type}}
\newcommand{\cells}[1]{\irs{cell}\left[#1\right]}
\newcommand{\vars}[1]{t_{#1}}

% Slot trees
\newcommand{\nilst}{\left[\right]}
\newcommand{\seqst}[1]{\left[#1\right]}
\newcommand{\tagst}[1]{\irs{tag}\left[#1\right]}
\newcommand{\flagsst}[3]{\irs{flags}\left<#1\right>\left[#2\tagsepst{}#3\right]}
\newcommand{\combst}[1]{\text{$<$}#1\text{$>$}}
\newcommand{\tagsepst}{\text{ | }}
\newcommand{\concatst}[2]{#1\oplus#2}

% Slot offset and data location
\newcommand{\offs}[2]{#1\text{$\sim$}#2}
\newcommand{\dataol}[2]{#1.#2}
\newcommand{\datal}[3]{\dataol{#1}{\offs{#2}{#3}}}

% Some notation commands
\newcommand{\sizeof}[1]{\phi\bigl(#1\bigr)}
\newcommand{\dual}[1]{\overline{#1}}
\newcommand{\getinstructions}[3]{\mathcal{I}_P(#1,\,#2,\,#3)}
\newcommand{\getendpoints}[2]{\mathcal{E}(#1,\,#2)}
\newcommand{\getexpression}[2]{\mathcal{I}_E(#1,\,#2)}
\newcommand{\makeblock}[1]{#1 \leftarrow \text{new}_{\text{b}}()}
\newcommand{\makecont}[1]{#1 \leftarrow \text{new}_{\text{c}}()}
\newcommand{\makedata}[2]{#1 \leftarrow \text{new}_{\text{d}}(#2)}
\newcommand{\makedrop}[3]{#1 \leftarrow \text{new}_{\text{e}}(#2,\,#3)}

% Type names
\newcommand{\classt}[1]{\textbf{\textcolor{forestgreen}{#1}}}
\newcommand{\closet}{\classt{close}}
\newcommand{\waitt}{\classt{wait}}
\newcommand{\lintt}{\classt{lint}}
\newcommand{\colintt}{\classt{colint}}
\newcommand{\lboolt}{\classt{lbool}}
\newcommand{\colboolt}{\classt{colbool}}
\newcommand{\lstringt}{\classt{lstring}}
\newcommand{\colstringt}{\classt{colstring}}
\newcommand{\choicet}[1]{\classt{choice of}\,\{#1\}}
\newcommand{\offert}[1]{\classt{offer of}\,\{#1\}}
\newcommand{\labelt}[2]{\classt{$|$\##1}:#2}
\newcommand{\sendt}[2]{\classt{send}\,#1;#2}
\newcommand{\recvt}[2]{\classt{recv}\,#1;#2}
\newcommand{\bangt}[1]{\classt{!}#1}
\newcommand{\whynott}[1]{\classt{?}#1}
\newcommand{\sendtyt}[2]{\classt{sendty}\,#1;#2}
\newcommand{\recvtyt}[2]{\classt{recvty}\,#1;#2}
\newcommand{\statet}[1]{\classt{state}\,#1}
\newcommand{\statelt}[1]{\classt{statel}\,#1}
\newcommand{\usaget}[1]{\classt{usage}\,#1}
\newcommand{\usagelt}[1]{\classt{usagel}\,#1}
\newcommand{\affinet}[1]{\classt{affine}\,#1}
\newcommand{\coaffinet}[1]{\classt{coaffine}\,#1}
\newcommand{\dualt}[1]{\classt{$\sim$}#1}
\newcommand{\rect}[2]{\classt{rec}\ #1;\,#2}
\newcommand{\corect}[2]{\classt{corec}\ #1;\,#2}

% Process names
\newcommand{\classp}[1]{\textbf{\textcolor{darkblue}{#1}}}
\newcommand{\emptyp}{\classp{()}}
\newcommand{\fwdp}{\classp{fwd}}
\newcommand{\cutp}{\classp{cut}}
\newcommand{\parp}{\classp{par}}
\newcommand{\closep}{\classp{close}}
\newcommand{\waitp}{\classp{wait}}
\newcommand{\casep}{\classp{case of}}
\newcommand{\labelp}[1]{\classp{\##1}}
\newcommand{\sendp}{\classp{send}}
\newcommand{\recvp}{\classp{recv}}
\newcommand{\bangp}{\classp{!}}
\newcommand{\whynotp}{\classp{?}}
\newcommand{\callp}{\classp{call}}
\newcommand{\sendtyp}{\classp{sendty}}
\newcommand{\recvtyp}{\classp{recvty}}
\newcommand{\unfoldmup}{\classp{unfold$_\mu$}}
\newcommand{\unfoldnup}{\classp{unfold$_\nu$}}
\newcommand{\corecp}{\classp{corec}}
\newcommand{\sharep}{\classp{share}}
\newcommand{\affinep}{\classp{affine}}
\newcommand{\discardp}{\classp{discard}}
\newcommand{\usep}{\classp{use}}
\newcommand{\cellp}{\classp{cell}}
\newcommand{\releasep}{\classp{release}}
\newcommand{\takep}{\classp{take}}
\newcommand{\putp}{\classp{put}}
\newcommand{\ifp}{\classp{if}}
\newcommand{\elsep}{\classp{else}}
\newcommand{\letp}{\classp{let}}
\newcommand{\scanp}{\classp{scan}}
\newcommand{\printp}{\classp{print}}
\newcommand{\printlnp}{\classp{println}}

% CLASS
\newcommand{\fwdc}[2]{\fwdp\,#1\,#2}
\newcommand{\cutc}[3]{\cutp\,\{#1\,\classp{$|$}#2\classp{$|$}\,#3\}}
\newcommand{\parc}[2]{\parp\,\{#1\,\classp{$||$}\,#2\}}
\newcommand{\emptyc}{\emptyp}
\newcommand{\closec}[1]{\closep\,#1}
\newcommand{\waitc}[1]{\waitp\,#1;}
\newcommand{\choosec}[2]{\classp{\##1}\ #2;}
\newcommand{\casec}[2]{\classp{case}\,#1\,\classp{of}\,\{#2\}}
\newcommand{\labelc}[2]{\classp{$|$}\labelp{#1}:#2}
\newcommand{\sendc}[3]{\sendp\,#1(#2.\,#3);}
\newcommand{\recvc}[2]{\recvp\,#1(#2);}
\newcommand{\bangc}[2]{\bangp\,#1(#2);}
\newcommand{\whynotc}[1]{\whynotp\,#1;}
\newcommand{\callc}[2]{\callp\,#1(#2);}
\newcommand{\sendtyc}[2]{\sendtyp\,#1(#2);}
\newcommand{\recvtyc}[2]{\recvtyp\,#1(#2);}
\newcommand{\unfoldmuc}[1]{\unfoldmup\,#1;}
\newcommand{\unfoldnuc}[1]{\unfoldnup\,#1;}
\newcommand{\corecc}[4]{\corecp\,#1(#2);\,#3\ [#4]}
\newcommand{\typedefc}[2]{\classp{type}\ #1\,\{#2\};;}
\newcommand{\rectypedefc}[2]{\classp{rec}\ \typedefc{#1}{#2}}
\newcommand{\typedefgc}[3]{\classp{type}\ #1(#2)\,\{#3\};;}
\newcommand{\rectypedefgc}[3]{\classp{rec}\ \typedefgc{#1}{#2}{#3}}
\newcommand{\procdefc}[3]{\classp{proc}\ \texttt{#1}(#2)\,\{#3\};;}
\newcommand{\procdefgc}[4]{\classp{proc}\ \texttt{#1}\langle#2\rangle(#3)\,\{#4\};;}
\newcommand{\recprocdefc}[3]{\classp{rec}\ \procdefc{#1}{#2}{#3}}
\newcommand{\recprocdefgc}[4]{\classp{rec}\ \procdefgc{#1}{#2}{#3}{#4}}
\newcommand{\proccallc}[2]{\texttt{#1}(#2)}
\newcommand{\proccallgc}[3]{\texttt{#1}\langle#2\rangle(#3)}
\newcommand{\sharec}[3]{\sharep\,#1\,\{#2\,\classp{$||$}\,#3\}}
\newcommand{\affinec}[1]{\affinep\,#1;}
\newcommand{\discardc}[1]{\discardp\,#1}
\newcommand{\usec}[1]{\usep\,#1;}
\newcommand{\cellc}[3]{\cellp\,#1(#2.\,#3)}
\newcommand{\cellfc}[2]{\cellp\,#1(#2)}
\newcommand{\releasec}[1]{\releasep\,#1}
\newcommand{\takec}[2]{\takep\,#1(#2);}
\newcommand{\putc}[3]{\putp\,#1(#2.\,#3);}
\newcommand{\putfc}[2]{\putp\,#1(#2);}
\newcommand{\ifc}[3]{\ifp\,#1\,\{#2\}\,\elsep\,\{#3\}}
\newcommand{\letc}[2]{\letp\,#1\,#2}
\newcommand{\scanc}[1]{\scanp(#1)}
\newcommand{\printc}[1]{\printp(#1);}
\newcommand{\printlnc}[1]{\printlnp(#1);}

% IR Data Locations
\newcommand{\slotoffset}[2]{#1\text{$\sim$}#2}
\newcommand{\offsetpast}[1]{\text{past}(#1)}
\newcommand{\offsetfuture}[1]{\text{future}(#1)}

% CLASS definitions
\newcommand{\slots}[2]{\mathcal{S}(#1,\,#2)}
\newcommand{\allslots}[2]{\mathcal{S}_T(#1,\,#2)}
\newcommand{\alllocalslots}[2]{\mathcal{S}_L(#1,\,#2)}

% Compiler environment accessors
\newcommand{\polarityof}[2]{#1_{p}(#2)}
\newcommand{\typeidof}[2]{#1_{t}(#2)}
\newcommand{\continuationof}[2]{#1_c(#2)}
\newcommand{\localdataof}[2]{#1_d(#2)}
\newcommand{\offsetof}[2]{#1_o(#2)}
\newcommand{\readlocof}[2]{#1_r(#2)}
\newcommand{\writelocof}[2]{#1_w(#2)}
\newcommand{\advancech}[6]{#1[#2 \overset{#3,#4}{\oplus} (#5,\,#6)]}

% C generator notation
\newcommand{\arch}{\mathcal{T}}
\newcommand{\archlayout}[1]{\arch(\texttt{#1})}
\newcommand{\slotlayout}[2]{\mathcal{L}_{\arch}(#1,\,#2)}
\newcommand{\layoutofvar}[2]{#1(#2)}
\newcommand{\envlayout}{\mathcal{E}}
\newcommand{\envlayouteps}[1]{\envlayout_\text{eps}(#1)}
\newcommand{\envlayoutdrop}[1]{\envlayout_\text{e}(#1)}
\newcommand{\envlayouttype}[2]{\envlayout_\text{ty}(#1, #2)}
\newcommand{\envlayoutcont}[2]{\envlayout_\text{c}(#1, #2)}
\newcommand{\envlayoutlocal}[2]{\envlayout_\text{d}(#1, #2)}

% Optimizer notation
\newcommand{\typeenv}{\theta}
\newcommand{\isvalue}[1]{\mathcal{V}_\typeenv(#1)}
\newcommand{\isdroppable}[1]{\mathcal{D}_\typeenv(#1)}
\newcommand{\iscloneable}[1]{\mathcal{C}_\typeenv(#1)}
\newcommand{\impossible}{\mathbf{impossible}}
\newcommand{\guaranteed}{\mathbf{guaranteed}}
\newcommand{\dependent}[1]{\mathbf{dependent}(#1)}
\newcommand{\depisvalue}[1]{\mathcal{V}(#1)}
\newcommand{\depisdroppable}[1]{\mathcal{D}(#1)}
\newcommand{\depiscloneable}[1]{\mathcal{C}(#1)}

% Session modes
\newcommand{\writemode}{\texttt{write}}
\newcommand{\readmode}{\texttt{read}}

% Indented CLASS environment
\newcounter{indentlevel}
\newlength{\indentstep}
\newenvironment{class}
{
    \begingroup

    \setcounter{indentlevel}{0}

    \setlength{\indentstep}{1.5em}

    \newcommand{\iz}{\addtocounter{indentlevel}{1}}
    \newcommand{\dz}{\addtocounter{indentlevel}{-1}}
    \newcommand{\pz}{\hspace*{\the\indentstep*\value{indentlevel}}}
    \newcommand{\z}{\\&\pz}

    % CLASS
    \let\cuti\cutc
    \renewcommand{\cutc}[3]{\cuti{\iz\z##1\z}{##2}{\z##3\dz\z}}
    \let\pari\parc
    \renewcommand{\parc}[2]{\pari{\iz\z##1\z}{\z##2\dz\z}}
    \let\waiti\waitc
    \renewcommand{\waitc}[1]{\waiti{##1}\z}
    \let\sendi\sendc
    \renewcommand{\sendc}[3]{\sendi{##1}{##2}{##3}\z}
    \let\recvi\recvc
    \renewcommand{\recvc}[2]{\recvi{##1}{##2}\z}
    \let\choosei\choosec
    \renewcommand{\choosec}[2]{\choosei{##1}{##2}\z}
    \let\casei\casec
    \renewcommand{\casec}[2]{\casei{##1}{##2\z}}
    \let\labeli\labelc
    \renewcommand{\labelc}[2]{\iz\z\labeli{##1}{##2}\dz}
    \let\unfoldmui\unfoldmuc
    \renewcommand{\unfoldmuc}[1]{\unfoldmui{##1}\z}
    \let\unfoldnui\unfoldnuc
    \renewcommand{\unfoldnuc}[1]{\unfoldnui{##1}\z}
    \let\procdefi\procdefc
    \renewcommand{\procdefc}[3]{\procdefi{##1}{##2}{\iz\z##3\dz\z}}
    \let\procdefgi\procdefgc
    \renewcommand{\procdefgc}[4]{\procdefgi{##1}{##2}{##3}{\iz\z##4\dz\z}}
    \let\typedefi\typedefc
    \renewcommand{\typedefc}[2]{\typedefi{##1}{\iz\z##2\dz\z}}
    \let\typedefgi\typedefgc
    \renewcommand{\typedefgc}[3]{\typedefgi{##1}{##2}{\iz\z##3\dz\z}}
    \let\printi\printc
    \renewcommand{\printc}[1]{\printi{##1}\z}
    \let\printlni\printlnc
    \renewcommand{\printlnc}[1]{\printlni{##1}\z}
    \let\affinei\affinec
    \renewcommand{\affinec}[1]{\affinei{##1}\z}
    \let\usei\usec
    \renewcommand{\usec}[1]{\usei{##1}\z}
    \let\discardi\discardc
    \renewcommand{\discardc}[1]{\discardi{##1}}
    \let\celli\cellc
    \renewcommand{\cellc}[3]{\celli{##1}{##2}{\iz\z##3\dz\z}}
    \let\cellfi\cellfc
    \renewcommand{\cellfc}[2]{\cellfi{##1}{##2}\z}
    \let\takei\takec
    \renewcommand{\takec}[2]{\takei{##1}{##2}\z}
    \let\puti\putc
    \renewcommand{\putc}[3]{\puti{##1}{##2}{##3}\z}
    \let\putfi\putfc
    \renewcommand{\putfc}[2]{\putfi{##1}{##2}\z}
    \let\sharei\sharec
    \renewcommand{\sharec}[3]{\sharei{##1}{\iz\z##2\z}{\z##3\dz\z}}

    $$
        \begin{aligned}
             & \pz
            }
            {
        \end{aligned}
    $$

    \endgroup
}

% Compilation environment
\newenvironment{compilation}[4]
{
    \begingroup
    % IR
    \newcommand{\beginir}{\\&\hspace{2.5em}}
    \newcommand{\mkzeroargir}[2]{\newcommand{##1}{\beginir\texttt{##2}()}}
    \newcommand{\mkoneargir}[2]{\newcommand{##1}[1]{\beginir\texttt{##2}(####1)}}
    \newcommand{\mktwoargir}[2]{\newcommand{##1}[2]{\beginir\texttt{##2}(####1, ####2)}}
    \newcommand{\mkthreeargir}[2]{\newcommand{##1}[3]{\beginir\texttt{##2}(####1, ####2, ####3)}}
    \newcommand{\mkfourargir}[2]{\newcommand{##1}[4]{\beginir\texttt{##2}(####1, ####2, ####3, ####4)}}
    \newcommand{\labelir}[1]{\\&\hspace{1.25em}\texttt{##1:}}
    \renewcommand{\comment}[1]{\quad\textcolor{forestgreen}{\%\ ##1}}
    \newcommand{\commentline}[1]{\beginir\textcolor{forestgreen}{\%\ ##1}}

    \mkoneargir{\launchthread}{launchThread}
    \mkoneargir{\pushtask}{pushTask}
    \mkoneargir{\poptask}{popTask}
    \mktwoargir{\finishcont}{finish}
    \mkoneargir{\continue}{continue}
    \mkthreeargir{\bindcontinuation}{bindContinuation}
    \mkthreeargir{\forward}{forward}
    \mktwoargir{\branchtag}{branchTag}
    \mktwoargir{\branchflags}{branchFlags}
    \mktwoargir{\writetype}{writeType}
    \mktwoargir{\writetag}{writeTag}
    \mktwoargir{\writeexpr}{writeExpression}
    \mktwoargir{\writescan}{writeScan}
    \mktwoargir{\writecell}{writeCell}
    \mkthreeargir{\writeexp}{writeExponential}
    \mktwoargir{\writecontinuation}{writeContinuation}
    \mkoneargir{\lockcell}{lockCell}
    \mkoneargir{\unlockcell}{unlockCell}
    \mkoneargir{\acquirecell}{acquireCell}
    \mkoneargir{\releasecell}{releaseCell}
    \mkoneargir{\print}{print}
    \mkoneargir{\printline}{printLine}
    \mkfourargir{\continueifneg}{continueIfNegative}
    \mkthreeargir{\moveslots}{moveSlots}
    \mkthreeargir{\cloneslots}{cloneSlots}
    \mktwoargir{\dropslots}{dropSlots}
    \mkthreeargir{\callprocess}{callProcess}
    \mkoneargir{\deferdrop}{deferDrop}

    \mkoneargir{\newtask}{newTask}
    \mkzeroargir{\nexttask}{nextTask}
    \mkthreeargir{\initcontinuation}{initializeContinuation}
    \mkoneargir{\freesession}{freeSession}
    \mkthreeargir{\call}{call}
    \mkoneargir{\flip}{flip}
    \mktwoargir{\fwd}{fwd}
    \mkoneargir{\pushclose}{pushClose}
    \mkoneargir{\popclose}{popClose}
    \mktwoargir{\pushtag}{pushTag}
    \mktwoargir{\poptag}{popTag}
    \mktwoargir{\pushsession}{pushSession}
    \mktwoargir{\popsession}{popSession}
    \newcommand{\hasexitpoint}{\texttt{exit point}}
    \newcommand{\recurse}[3]{\beginir\getinstructions{##1}{##2}{##3}}

    $$
        \begin{aligned}
             & \getendpoints{#1}{#3} \equiv #4 \\
             & \getinstructions{#1}{#2}{#3} \equiv
            }
            {
        \end{aligned}
    $$
    \endgroup
}

% For \Mapsto
\usepackage{stmaryrd}

% For centering tables
\usepackage{tabularx}
